\documentclass[12pt,a4paper]{article}
\usepackage[utf8]{inputenc}
\usepackage[T1]{fontenc}
\usepackage[spanish]{babel}
\usepackage{amsmath}
\usepackage{amssymb}
\usepackage{geometry}
\usepackage{hyperref}
\usepackage{bookmark}
\usepackage{fancyhdr}
\usepackage{enumitem}
\usepackage{booktabs}
\usepackage{titlesec}
\geometry{margin=2.5cm}

% Encabezado estilo notas de clase
\setlength{\headheight}{15pt}
\pagestyle{fancy}
\fancyhf{}
\fancyhead[L]{Microestructura, Eficiencia, Renta Fija y Derivados}
\fancyhead[R]{Problem Set \#4}
\fancyfoot[C]{\thepage}
\fancyfoot[L]{TA Session \#4}
\fancyfoot[R]{Jesus Villota Miranda}
\renewcommand{\footrulewidth}{0.4pt}

% Encabezado
\title{\textbf{Problem Set 4}\\\large Microestructura, Eficiencia, Renta Fija y Derivados}
\author{Jesus Villota Miranda\\CEMFI}
\date{\today}

\newcommand{\E}{\mathbb{E}}
\newcommand{\V}{\mathbb{V}}
\newcommand{\Cov}{\mathbb{C}\text{ov}}
\newcommand{\R}{\mathbb{R}}
\newcommand{\EUR}[1]{\ensuremath{#1~\text{EUR}}}

% Formato de ejercicios
\newenvironment{ejercicio}[1]{\subsection*{Ejercicio #1}}{}

\begin{document}

% \maketitle

\section{Microestructura y negociación}

\begin{ejercicio}{1 (Compra al margen y llamada de mantenimiento)}
Una inversora compra \(400\) acciones de una compañía a \EUR{35} cada una financiando \EUR{4{,}500} con su bróker. El préstamo devenga un interés anual del \(9\%\).
\begin{enumerate}[label=\alph*)]
  \item Calcule el margen inicial en euros y como porcentaje del valor de la posición.
  \item Si al cabo de un año el precio cae a \EUR{28} por acción, determine el valor del margen tras pagar intereses.
  \item Con un margen de mantenimiento del \(30\%\), ¿hay llamada de margen? Justifique con el porcentaje de margen.
  \item Calcule la rentabilidad de la inversora sobre su capital propio y compárela con la variación porcentual del precio.
\end{enumerate}
\end{ejercicio}

\begin{ejercicio}{2 (Venta corta con dividendo y mantenimiento)}
Un inversor abre una posición corta de \(800\) acciones a \EUR{48} con margen inicial del \(50\%\). La acción paga un dividendo de \EUR{1{,}50} dentro del año. Al cierre del año el precio sube a \EUR{55}.
\begin{enumerate}[label=\alph*)]
  \item Calcule el efectivo obtenido y el margen inicial en euros.
  \item Determine el margen restante después del dividendo y del cambio de precio.
  \item Si el margen de mantenimiento es \(35\%\), indique si se produce llamada de margen.
  \item Compute la rentabilidad de la posición corta tomando el margen inicial como inversión.
\end{enumerate}
\end{ejercicio}

\begin{ejercicio}{3 (Libro de órdenes y costes de transacción)}
Considere el siguiente libro de órdenes limitadas para la acción Soltec. La última operación fue a \EUR{62}.
\begin{center}
\begin{tabular}{cc|cc}
\toprule
\multicolumn{2}{c|}{Órdenes de compra} & \multicolumn{2}{c}{Órdenes de venta} \\
Precio & Cantidad & Precio & Cantidad\\
\midrule
61{,}80 & 400 & 62{,}30 & 200\\
61{,}50 & 600 & 63{,}10 & 300\\
61{,}20 & 500 & 64{,}00 & 400\\
60{,}80 & 900 & 65{,}50 & 600\\
\bottomrule
\end{tabular}
\end{center}
\begin{enumerate}[label=\alph*)]
  \item Determine el \emph{bid-ask spread} y el tamaño medio a ambos lados en el mejor nivel.
  \item Una orden de compra a mercado por \(300\) acciones, ¿a qué precio medio efectivo se ejecuta?
  \item ¿Cuál sería el precio medio y el coste total si la orden de compra fuera de \(900\) acciones?
  \item Calcule el coste de transacción implícito (en euros y como \% del valor nocional) de la orden de \(900\) acciones usando el último precio de \EUR{62} como referencia.
\end{enumerate}
\end{ejercicio}

\section{Eficiencia de mercado y eventos}

\begin{ejercicio}{4 (Rentabilidad anormal y CAR con CAPM)}
Una acción tiene \(\alpha=0\) y \(\beta=1{,}2\) estimados de forma estable. Durante tres meses consecutivos los rendimientos mensuales del mercado fueron \(3\%\), \(1\%\) y \(-2\%\). Los rendimientos observados de la acción fueron \(4{,}6\%\), \(2{,}0\%\) y \(-0{,}5\%\). Suponga \(r_f=0\).
\begin{enumerate}[label=\alph*)]
  \item Calcule el rendimiento esperado en cada mes según el CAPM.
  \item Obtenga los rendimientos anormales mensuales.
  \item Compute el CAR (cumulative abnormal return) del trimestre y evalúe si indica noticia positiva o negativa.
\end{enumerate}
\end{ejercicio}

\section{Renta fija}

\begin{ejercicio}{6 (Comparación de bonos: cupón cero vs cupón anual)}
Considere dos bonos con nominal \EUR{1{,}000} y vencimiento a 2 años:
\begin{itemize}
   \item \textbf{Bono A:} Bono cupón cero (zero coupon bond)
   \item \textbf{Bono B:} Bono con cupón anual del \(5\%\)
\end{itemize}
Los tipos spot anuales son: \(r_1=3{,}0\%\) y \(r_2=3{,}5\%\). Los bonos fueron adquiridos a los siguientes precios de compra: Bono A a \EUR{925} y Bono B a \EUR{1{,}035}.
\begin{enumerate}[label=\alph*)]
   \item Calcule el precio teórico de ambos bonos usando la curva spot.
  \item Asumiendo que los bonos se mantienen hasta el vencimiento, calcule la TIR (yield to maturity) de cada bono. La TIR es la tasa única de descuento que iguala el precio de compra con el valor presente de los flujos futuros. Para el Bono B, resuelva la ecuación cuadrática correspondiente para calcular la TIR. Comente la relación entre la TIR y los tipos spot, y explique por qué la TIR del Bono A difiere del tipo spot \(r_2\) cuando el precio de compra no coincide con el teórico.
   \item Calcule la duración de Macaulay de ambos bonos usando sus respectivas TIR como tasa de descuento uniforme.
   \item La nueva curva de tipos muestra que \(r_1 = 3{,}5\%\) y el tipo forward implícito \(f_{1,1} = 4{,}5\%\). Calcule primero el tipo spot \(r_2\) a partir de \(r_1\) y \(f_{1,1}\). Luego, calcule el nuevo precio teórico de ambos bonos usando la nueva curva spot.
\end{enumerate}
\end{ejercicio}

\section{Derivados: forwards}

\begin{ejercicio}{8 (Curva spot y tipos forward implícitos)}
Se observa la siguiente curva de tipos spot anuales (compuestos anualmente):
\begin{center}
\begin{tabular}{cc}
\toprule
Vencimiento (años) & Tipo spot anual (\%) \\
\midrule
1 & \(3{,}0\) \\
2 & \(3{,}5\) \\
3 & \(4{,}0\) \\
4 & \(4{,}2\) \\
5 & \(4{,}3\) \\
\bottomrule
\end{tabular}
\end{center}
\begin{enumerate}[label=\alph*)]
   \item Calcule el tipo forward implícito \(f_{1,1}\) (forward de 1 año que comienza dentro de 1 año).
   \item Calcule el tipo forward implícito \(f_{2,1}\) (forward de 1 año que comienza dentro de 2 años).
   \item Calcule el tipo forward implícito \(f_{1,2}\) (forward de 2 años que comienza dentro de 1 año).
   \item Si un inversor quiere asegurar hoy un tipo de interés del \(4{,}5\%\) para un préstamo de 1 año que comenzará dentro de 2 años, ¿le conviene más contratar un forward o esperar? Justifique su respuesta comparando con el tipo forward implícito.
   \item Explique intuitivamente por qué el tipo forward \(f_{1,1}\) es mayor que el tipo spot a 1 año \(r_1\) cuando la curva tiene pendiente positiva.
\end{enumerate}
\end{ejercicio}

\begin{ejercicio}{9 (Libro de órdenes y bonos corporativos AirBaltic)}
El bono corporativo de AirBaltic fue emitido el 15 de noviembre de 2024 con nominal de \EUR{1{,}000} cada uno, paga cupones trimestrales a una tasa anual del \(14{,}5\%\) y madura el 15 de noviembre de 2028. El bono es callable. El libro de órdenes muestra las siguientes posiciones (precios en \% del nominal, sin cupón corrido):
\begin{center}
\begin{tabular}{cc|cc}
\toprule
\multicolumn{2}{c|}{Órdenes de compra} & \multicolumn{2}{c}{Órdenes de venta} \\
Precio & Cantidad (nom.) & Precio & Cantidad (nom.)\\
\midrule
100{,}40 & 200{,}000 & 101{,}00 & 150{,}000\\
100{,}25 & 250{,}000 & 101{,}20 & 150{,}000\\
99{,}90  & 300{,}000 & 101{,}50 & 200{,}000\\
         &           & 102{,}00 & 100{,}000\\
\bottomrule
\end{tabular}
\end{center}
\begin{enumerate}[label=\alph*)]
  \item El 15 de diciembre de 2025, Jesus envía una orden de compra a mercado por nominal \EUR{350{,}000}. ¿A qué precio efectivo (precio medio ponderado en \% del nominal) adquiere el bono? Calcule también el cupón corrido que debe pagar en la fecha de compra y el desembolso total en euros.
  \item Si Jesus mantuviera el bono hasta su vencimiento (15 de noviembre de 2028), plantee la ecuación que debería satisfacer la TIR (YTM) anual del bono. Considere que el bono paga cupones trimestrales y devuelve el \(100\%\) del nominal al vencimiento, y que dado que el bono se compró el 15 de diciembre de 2025, los primeros pagos no coinciden con trimestres completos desde la compra.
  \item Sin embargo, Jesus decide vender el bono el 15 de abril de 2026. En esa fecha, el libro de órdenes muestra las siguientes posiciones (precios en \% del nominal, sin cupón corrido):
  \begin{center}
  \begin{tabular}{cc|cc}
  \toprule
  \multicolumn{2}{c|}{Órdenes de compra} & \multicolumn{2}{c}{Órdenes de venta} \\
  Precio & Cantidad (nom.) & Precio & Cantidad (nom.)\\
  \midrule
  99{,}20 & 100{,}000 & 99{,}80 & 200{,}000\\
  99{,}00 & 150{,}000 & 100{,}00 & 250{,}000\\
  98{,}75 & 200{,}000 & 100{,}30 & 300{,}000\\
  98{,}40 & 150{,}000 &         &          \\
  \bottomrule
  \end{tabular}
  \end{center}
  Calcule el precio efectivo de venta que obtiene Jesus al vender los \EUR{350{,}000} nominales mediante una orden de venta a mercado. Luego, determine la rentabilidad anual (TIR) obtenida por Jesus en esta inversión, considerando que al vender en abril tiene derecho a cobrar la parte proporcional del cupón trimestral (2/3 del cupón, ya que faltan 30 días para el pago del 15 de mayo).
\end{enumerate}
\end{ejercicio}

\section{Derivados: forwards, futures y opciones}

\begin{ejercicio}{10 (Valoración de forwards y payoffs de opciones)}
Una acción de la empresa TechCorp cotiza hoy a \EUR{45} y no paga dividendos. El tipo de interés libre de riesgo anual es del \(4\%\).
\begin{enumerate}[label=\alph*)]
  \item Calcule el precio forward teórico (precio de ejercicio \(K\)) para un contrato forward que vence dentro de 2 años usando el principio de no arbitraje.
  \item Si hace 6 meses se firmó un contrato forward para comprar la acción a \EUR{48} dentro de 2 años, y hoy la acción cotiza a \EUR{50}, calcule el valor actual del contrato forward para el comprador (posición larga).
  \item Un inversor compra una opción de compra (call) europea sobre TechCorp con precio de ejercicio \EUR{48} y vencimiento en 1 año, pagando una prima de \EUR{3{,}50}. Calcule el beneficio neto del inversor si al vencimiento el precio de la acción es: (i) \EUR{52}, (ii) \EUR{48}, (iii) \EUR{44}.
  \item El mismo inversor compra también una opción de venta (put) europea sobre TechCorp con precio de ejercicio \EUR{48} y vencimiento en 1 año, pagando una prima de \EUR{2{,}00}. Calcule el beneficio neto total si al vencimiento el precio de la acción es \EUR{44} (considerando ambas opciones).
\end{enumerate}
\end{ejercicio}

\end{document}

